\documentclass[UTF8,fontset=fandol,10pt]{ctexart}
\usepackage[a4paper,margin=18mm,top=20mm,bottom=19mm]{geometry}
\usepackage[most]{tcolorbox}
\usepackage{enumitem,fancyhdr}
\definecolor{Navy}{HTML}{245781}
\definecolor{Risk}{HTML}{AC161C}
\definecolor{Good}{HTML}{167333}
\definecolor{Gold}{HTML}{9B8B39}
\definecolor{Slate}{HTML}{526273}
\pagestyle{fancy}\fancyhf{}
\setlength{\headheight}{15pt}
\fancyfoot[C]{\small\thepage}
\setlength{\parindent}{0pt}\setlength{\parskip}{4pt}
\setlist[itemize]{leftmargin=1.2em,itemsep=2pt,topsep=1pt}
\ctexset{section={format=\large\bfseries,beforeskip=12pt,afterskip=7pt},subsection={format=\normalsize\bfseries}}
\raggedbottom
\fancyhead[L]{\small 虚构案例：内容运营面试复盘}\fancyhead[R]{\small 面试复盘}
\begin{document}\begin{center}{\large 面试复盘}\par\vspace{4mm}{\LARGE\bfseries 虚构案例：内容运营面试复盘}\par\vspace{3mm}{\small 虚构演示 · 内容运营 · 无时间戳文字稿}\end{center}\begin{tcolorbox}[enhanced,breakable,sharp corners,boxrule=.6pt,colback=Navy!5!white,colframe=Navy,colbacktitle=Navy,coltitle=white,fonttitle=\bfseries,title={本场结论与证据边界},left=2mm,right=2mm,top=1.5mm,bottom=1.5mm]个人职责边界清楚，具有初步比较意识；下一场重点补齐目标、判断依据和验证步骤。协作题因回答缺失，暂不评分。\par\small 全部材料均为虚构演示。三类材料齐全，但无时间戳，第三题回答缺失；仅评价记录支持的内容，不预测录用结果。\end{tcolorbox}\begin{minipage}[t]{.485\linewidth}\small \begin{tcolorbox}[enhanced,breakable,sharp corners,boxrule=.6pt,colback=Risk!5!white,colframe=Risk,colbacktitle=Risk,coltitle=white,fonttitle=\bfseries,title={Top 5 扣分点},left=2mm,right=2mm,top=1.5mm,bottom=1.5mm]\begin{itemize}\item q01未说明活动目标。\item q01未给出个人判断依据。\item q01缺少结果验证方式；没有数据本身不等同于表现差。\item q02未明确观察窗口及指标口径。\item 证据不足：q03待核，需补充实际回答。\end{itemize}\end{tcolorbox}\end{minipage}\hfill \begin{minipage}[t]{.485\linewidth}\small \begin{tcolorbox}[enhanced,breakable,sharp corners,boxrule=.6pt,colback=Good!5!white,colframe=Good,colbacktitle=Good,coltitle=white,fonttitle=\bfseries,title={Top 5 可复用亮点},left=2mm,right=2mm,top=1.5mm,bottom=1.5mm]\begin{itemize}\item q01明确个人执行职责。\item q01未夸大为独立主导。\item q01承认数据缺失。\item q02提出同账号比较。\item q02识别题材与发布时间差异。\end{itemize}\end{tcolorbox}\end{minipage}\par\vspace{2mm}\begin{minipage}[t]{.485\linewidth}\small \begin{tcolorbox}[enhanced,breakable,sharp corners,boxrule=.6pt,colback=Gold!5!white,colframe=Gold,colbacktitle=Gold,coltitle=white,fonttitle=\bfseries,title={需要改进的方向},left=2mm,right=2mm,top=1.5mm,bottom=1.5mm]\begin{itemize}\item q01补充活动目标及背景。\item q01说明本人判断和动作。\item q02先确认指标口径及观察窗口。\item q03补充真实分歧案例。\end{itemize}\end{tcolorbox}\end{minipage}\hfill \begin{minipage}[t]{.485\linewidth}\small \begin{tcolorbox}[enhanced,breakable,sharp corners,boxrule=.6pt,colback=Slate!5!white,colframe=Slate,colbacktitle=Slate,coltitle=white,fonttitle=\bfseries,title={下一场优先级},left=2mm,right=2mm,top=1.5mm,bottom=1.5mm]\begin{itemize}\item 优先练习q01的个人贡献边界。\item 优先练习q02的排查顺序。\item 补齐q03回答证据。\item 核对简历与练习稿的事实一致性。\end{itemize}\end{tcolorbox}\end{minipage}\par\vspace{2mm}\clearpage\section*{面试官问题与现场回答复盘}\begin{tcolorbox}[enhanced,breakable,sharp corners,boxrule=.6pt,colback=Navy!5!white,colframe=Navy,colbacktitle=Navy,coltitle=white,fonttitle=\bfseries,title={q01 · 请介绍你的活动经历。},left=2mm,right=2mm,top=1.5mm,bottom=1.5mm]\textbf{来源：}transcript:P1—P2\quad\textbf{评分：}3/5\par \textbf{具体问题：}请介绍你的活动经历。\par \textbf{现场回答（整理版）：}我参与过一次内容活动，主要整理素材、跟进发布。具体业绩我没有数据。\par \textbf{具体优点：}明确区分参与及个人执行内容，承认没有业绩数据。\par \textbf{具体失误：}没有讲清活动目标、个人判断和验证方式。\par \textbf{修改策略：}补充本人能确认的目标与动作；没有业绩数据时保留说明。\par \textbf{练习逐字稿：}我参与过一次内容活动，负责素材整理和发布跟进。活动目标是[待补充信息]。我采取的具体动作是[待补充信息]。目前没有可核对的个人业绩数据，不能把整体结果归为我的贡献。\par \end{tcolorbox}\begin{tcolorbox}[enhanced,breakable,sharp corners,boxrule=.6pt,colback=Navy!5!white,colframe=Navy,colbacktitle=Navy,coltitle=white,fonttitle=\bfseries,title={q02 · 如果播放量下降，你会如何排查？},left=2mm,right=2mm,top=1.5mm,bottom=1.5mm]\textbf{来源：}transcript:P3—P4\quad\textbf{评分：}3/5\par \textbf{具体问题：}如果播放量下降，你会如何排查？\par \textbf{现场回答（整理版）：}我会先比较同一账号最近几条内容，再看是不是题材或者发布时间不一样。\par \textbf{具体优点：}提出同账号比较，注意题材和发布时间差异。\par \textbf{具体失误：}尚未说明播放下降的口径、观察窗口以及如何验证假设。\par \textbf{修改策略：}先核对指标和窗口，再拆分可能因素，选一个可控因素验证。\par \textbf{练习逐字稿：}我会先确认播放量口径和观察窗口一致，再比较同账号内容，检查题材、发布时间及流量来源。对发现的差异先当作假设，选择一个因素做小范围对照，再看同一窗口内的结果。\par \end{tcolorbox}\begin{tcolorbox}[enhanced,breakable,sharp corners,boxrule=.6pt,colback=Navy!5!white,colframe=Navy,colbacktitle=Navy,coltitle=white,fonttitle=\bfseries,title={q03 · 你怎样处理和设计同事的分歧？},left=2mm,right=2mm,top=1.5mm,bottom=1.5mm]\textbf{来源：}transcript:P5—P6\quad\textbf{评分：}未评分：证据不足\par \textbf{具体问题：}你怎样处理和设计同事的分歧？\par \textbf{现场回答（整理版）：}证据不足：回答片段缺失。\par \textbf{具体优点：}本场未充分验证。\par \textbf{具体失误：}证据不足，不能判断是否存在失误。\par \textbf{修改策略：}补充原始回答后再评价，不从简历推测现场表现。\par \textbf{练习逐字稿：}待补充信息：请先回忆实际分歧、各方目标、个人动作和最终处理结果。\par \end{tcolorbox}\section*{动态重点追问}重点追问q01：哪些动作由你负责？依据来自哪里？重点追问q02：怎样排除观察窗口不同的影响？q03先补充证据。\par \section*{整体复盘}语言描述：q01职责明确，但目标和结果主线尚不完整。经历阐述：个人边界清楚，验证环节待补充。业务思考：q02具备比较意识，缺少口径与验证步骤。协作能力本场未充分验证。\par \section*{后续面试问题}\begin{itemize}\item 活动目标是什么？\item 你的个人职责是什么？\item 素材选择的判断依据是什么？\item 如何验证发布动作有效？\item 播放量比较的观察窗口如何统一？\item 如何区分内容和流量来源的影响？\item 如何设计单因素验证？\item 请补充一次真实协作分歧及处理过程。\end{itemize}\section*{巩固资料}\begin{itemize}\item 待补充资料1：本示例未核验公开链接；请针对指标定义查找官方资料并说明练习用途。\item 待补充资料2：本示例未核验公开链接；请针对对照试验查找官方资料并说明练习用途。\item 待补充资料3：本示例未核验公开链接；请针对活动复盘查找官方资料并说明练习用途。\item 待补充资料4：本示例未核验公开链接；请针对内容分析查找官方资料并说明练习用途。\item 待补充资料5：本示例未核验公开链接；请针对跨团队沟通查找官方资料并说明练习用途。\end{itemize}\end{document}